% ============================================================================
% APPENDIX A: REPRODUCIBILITY DETAILS
% ============================================================================
\appendix
\section{Reproducibility Details}\label{app:reproducibility}

This appendix records the computational protocol used for the numerical claims
in Section~\ref{sec:numerical}.  The computations are intentionally small and
exhaustive: for $s=4$, all input values, all key hypotheses, all Walsh
coordinates, and all active coordinate planes are enumerated.

\subsection{Software environment}

The distributed scripts use Python 3 and standard numerical
packages.  The core dependencies are \texttt{numpy} for exact finite-array
enumeration and linear algebra, and \texttt{matplotlib} for figure generation.
No randomized optimization, training procedure, or stochastic sampling is used
for the reported tables.  When a random 4-bit bijection is used in
\texttt{verify\_curvature.py}, the script fixes the pseudorandom seed explicitly.

\subsection{Commands}

The main numerical outputs can be regenerated with the following commands from
the project root:
\begin{verbatim}
python gen_figures.py
python compute_curvatures.py
python verify_curvature.py
\end{verbatim}
The first command regenerates the Walsh-spectrum and local key-separation
figures.  The second command constructs the boundary covariance matrix,
projects onto the active eigenspace, computes the projected third cumulant
contractions, and evaluates the curvature diagnostic.  The third command is a
sanity check for the value $K=-1/4$ on PRESENT and on an additional fixed
random bijection.

\subsection{Arithmetic and precision}

The Walsh coefficients, boundary covariance entries, and third cumulant entries
are obtained from finite sums over $\{0,1\}^4$.  These quantities are therefore
integer or rational before diagonalization.  The eigenspace projection and the
curvature table are evaluated in floating-point arithmetic after the exact
finite sums have been constructed.  The reported curvature values are therefore
reported as reproducible numerical diagnostics for the tested S-boxes.

\subsection{Sign conventions}

The Walsh transform convention is
\[
  \widehat S(\alpha,\beta)=\sum_{a\in\{0,1\}^s}
  (-1)^{\alpha\cdot a\oplus\beta\cdot S(a)}.
\]
The active-eigenspace curvature diagnostic uses the convention in
\eqref{eq:curvature-diagnostic}.  Changing the Riemann tensor sign convention
would change the sign of all reported sectional curvature values; the scripts
and the paper use the same convention.

\subsection{Reproducibility coverage}

The computational protocol covers the quantities reported in the paper: Walsh
spectra, boundary covariance matrices, active ranks/eigenvalues, local spectral
key-separation proxies, Hamming-weight leakage descriptors, and the projected
active-eigenspace curvature diagnostic for the tested 4-bit S-boxes.

\section*{Declarations}

\subsection*{Availability of data and materials}

The datasets supporting the conclusions of this article are included within the
article and its additional file.  The additional file
\texttt{supplementary\_material.zip} contains the scripts used to regenerate
the numerical tables and figures.

Software information: project name, Boundary Information Geometry for S-box
Side-Channel Leakage; project home page,
\url{https://github.com/EkodeckStephane/IG-based-Sbox-SCL}; archived version,
not applicable before repository release or DOI assignment;
operating system, platform independent; programming language, Python 3; other
requirements, \texttt{numpy} and \texttt{matplotlib}; license, not specified;
restrictions to use by non-academics, none beyond the repository license status.

\subsection*{Competing interests}

The authors declare that they have no competing interests.

\subsection*{Funding}

No specific funding was received for this work.

\subsection*{Authors' contributions}

Conceptualization and methodology: S.G.R.E. and S.A.E.; formal analysis:
S.G.R.E., S.A.E., P.R.M.A., and R.N.; software and reproducible computations:
S.G.R.E.; validation: S.A.E., P.R.M.A., and R.N.; supervision and project
administration: R.N.; writing--original draft: S.G.R.E.; writing--review and
editing: S.G.R.E., S.A.E., P.R.M.A., and R.N.  All authors read and approved
the final manuscript.

\subsection*{Acknowledgements}

Not applicable.

\subsection*{Authors' information}

Not applicable.
